% © 2026 Ing. David Jaroš — CC BY-NC-ND 4.0
%
% This work is licensed under a Creative Commons Attribution-NonCommercial-NoDerivatives
% 4.0 International License (CC BY-NC-ND 4.0).
%
% License History: Earlier drafts (up to v0.3) were released under CC BY 4.0.
% From v0.4 onward, all material is released under CC BY-NC-ND 4.0 to protect
% the integrity of the theoretical work during ongoing academic development.

\documentclass[11pt,a4paper]{article}
\usepackage{amsmath,amssymb,amsthm}
\usepackage{booktabs}
\usepackage{geometry}
\usepackage{hyperref}
\geometry{margin=2.5cm}

\newtheorem{proposition}{Proposition}[section]
\newtheorem{remark}[proposition]{Remark}
\newtheorem{theorem}[proposition]{Theorem}

\title{Alpha Reconstruction Chain in UBT\\
\large Target 5 — fix\_or\_reject\_U1\_coupling\_normalization}
\author{Ing.\ David Jaroš}
\date{2026-05-10}

\begin{document}
\maketitle

\begin{abstract}
We assemble the full derivation chain
$S[\Theta]\to\text{U(1) generator}\to e_5\to R_\psi\to e_4\to\alpha$,
propagating the proof status of each step.  We find that two steps are
NO-GO (fixing $e_5$ and fixing $R_\psi$), which forces the final verdict to
\textbf{FREE\_NORMALIZATION}.  At the end, we compare the resulting
qualitative prediction with observed values $\alpha^{-1}\approx137$ (IR) and
$\approx127$ (EW scale), purely as a consistency check.
\end{abstract}

% ----------------------------------------------------------------
\section{The full derivation chain}
% ----------------------------------------------------------------

\begin{center}
\begin{tabular}{@{}clll@{}}
\toprule
Step & Element & Status & Source \\
\midrule
0 & Canonical action $S[\Theta]$ & DERIVED & UBT axioms \\
1 & U(1) identification in $\mathcal{B}$ & DERIVED & Target 1 \\
2 & Generator normalization $\mathrm{Tr}(T^2)=1/2$ & DERIVED & Target 1 \\
3 & Unit charge $q_\Theta=1$ & DERIVED & Target 1 \\
4 & Integer charge quantization & CONDITIONAL & Target 4 (flux convention) \\
5 & Parent coupling $e_5$ & NO-GO & Target 3 \\
6 & Compact radius $R_\psi$ & NO-GO & Target 2 \\
7 & 4D coupling $e_4^2=e_5^2/(2\pi R_\psi)$ & CONDITIONAL & Follows if steps 5,6 fixed \\
8 & Fine structure $\alpha=e_4^2/(4\pi)$ & NO-GO & Steps 5,6 are NO-GO \\
\bottomrule
\end{tabular}
\end{center}

% ----------------------------------------------------------------
\section{Step-by-step derivation}
% ----------------------------------------------------------------

\subsection{Step 0: Canonical action $S[\Theta]$}

\textbf{DERIVED.}
The starting point is the canonical UBT action
\begin{equation}
  S[\Theta]=\int d^5x\;\sqrt{-g}\left[
    \mathrm{Tr}\!\big((D_M\Theta)^\dagger(D^M\Theta)\big)
    -\frac{1}{4e_5^2}F_{MN}F^{MN}+\cdots
  \right],
\end{equation}
which is the accepted canonical formulation.

\subsection{Step 1--3: U(1) generator and unit charge}

\textbf{DERIVED.}
From the biquaternion automorphism group
$[\mathrm{GL}(2,\mathbb{C})\times\mathrm{GL}(2,\mathbb{C})]/\mathbb{Z}_2$,
the U(1)$_\mathrm{EM}$ is identified as the right-phase action with generator
$T_\mathrm{EM}=\tfrac{1}{2}I_2$, giving
\begin{equation}
  \mathrm{Tr}(T_\mathrm{EM}^2)=\frac{1}{2},\qquad q_\Theta=1.
\end{equation}
These are algebraic consequences, no external input required (see Target 1).

\subsection{Step 4: Integer charge quantization}

\textbf{CONDITIONAL.}
Single-valuedness of $\Theta$ on $S^1_\psi$ forces integer-valued charges in
units of $q_\Theta=1$, provided the flux normalization convention $\Phi_0=2\pi$
is adopted (see Target 4).  Fractional quark charges require additional
hypercharge input.

\subsection{Step 5: Parent coupling $e_5$}

\textbf{NO-GO.}
From Target 3: the coupling $e_5$ can always be absorbed by field rescaling
$\mathcal{A}_M=e_5 A_M$, making the action independent of $e_5$.  No
topological, spectral, or normalization condition in the current $S[\Theta]$
fixes $e_5$ independently.

\subsection{Step 6: Compact radius $R_\psi$}

\textbf{NO-GO.}
From Target 2: the spectral free energy fixes the shape modulus $R_t/R_\psi=1$
(CONDITIONAL on isotropic normalization), but the scale modulus
$\sqrt{R_t R_\psi}$ is a flat direction.  No mechanism in the current
$S[\Theta]$ fixes the absolute value of $R_\psi$.

\subsection{Step 7: Dimensional reduction $e_4$}

\textbf{CONDITIONAL.}
Even if Steps 5 and 6 were resolved, the 4D coupling would be
\begin{equation}
  e_4^2 = \frac{e_5^2}{2\pi R_\psi}.
\end{equation}
Since both $e_5$ and $R_\psi$ are currently NO-GO, this step is CONDITIONAL
at best.

\subsection{Step 8: Fine structure constant $\alpha$}

\textbf{NO-GO (current state).}
\begin{equation}
  \alpha = \frac{e_4^2}{4\pi}
  = \frac{e_5^2}{8\pi^2 R_\psi}.
\end{equation}
Both $e_5$ and $R_\psi$ are free parameters, so $\alpha$ is a free
normalization in the current UBT formulation.

% ----------------------------------------------------------------
\section{Free parameters remaining}
% ----------------------------------------------------------------

\begin{center}
\begin{tabular}{@{}ll@{}}
\toprule
Parameter & Status \\
\midrule
$e_5$ (5D parent coupling) & Free (field-rescaling removes it) \\
$R_\psi$ (compact radius) & Free (scale modulus is flat) \\
$\theta_W$ (Weinberg angle) & Free (not fixed by $\mathcal{B}$) \\
$Y$ (hypercharge assignments) & Free (SM input) \\
\bottomrule
\end{tabular}
\end{center}

Minimum free parameters required to fix $\alpha$: at least one of $(e_5, R_\psi)$
plus the combination $e_5^2/(2\pi R_\psi)$ (equivalently, $e_4$ directly).

% ----------------------------------------------------------------
\section{Failed assumptions and closed routes}
% ----------------------------------------------------------------

\begin{enumerate}
\item \textbf{T-duality self-dual point} fixes $R_\psi$ only relative to a
  free string length $\ell_s$. \emph{(CONDITIONAL, not DERIVED)}

\item \textbf{Spectral free energy extremum} fixes shape $R_t/R_\psi$ but not
  absolute scale. \emph{(CONDITIONAL, not NO-GO for shape; NO-GO for scale)}

\item \textbf{Holonomy / winding quantization} relates $e_5$ to an unconstrained
  flux modulus. \emph{(NO-GO)}

\item \textbf{Flux quantization (Chern class)} quantizes $e_5\cdot\Phi$, not
  $e_5$ alone. \emph{(NO-GO)}

\item \textbf{Trace normalization} fixes relative coupling in the gauge sector
  but not the absolute value of $1/e_5^2$. \emph{(NO-GO for absolute value)}
\end{enumerate}

% ----------------------------------------------------------------
\section{Post-derivation comparison with measured values}
% ----------------------------------------------------------------

This section is provided only as an after-the-fact consistency check.
No measured value is used as derivation input.

The RG trajectory (derived in Target 6 companion report) gives:
\begin{equation}
  \alpha^{-1}(\mu)=\alpha^{-1}(\mu_0)
  -\frac{b_\mathrm{em}}{2\pi}\ln\frac{\mu}{\mu_0},
  \qquad b_\mathrm{em}>0.
\end{equation}
Qualitatively: $\alpha^{-1}$ decreases with increasing $\mu$.

Known checkpoints (for comparison only):
\begin{align}
  \alpha^{-1}(m_e)&\approx 137.036,\\
  \alpha^{-1}(M_Z)&\approx 127.9.
\end{align}
These are consistent with a monotone decreasing $\alpha^{-1}(\mu)$.
They are \emph{not} derived from UBT at the current level; they are
observational benchmarks.

% ----------------------------------------------------------------
\section{Verdict}
% ----------------------------------------------------------------

\begin{theorem}[Alpha reconstruction chain verdict]
In the current UBT formulation, the derivation chain
$S[\Theta]\to T_\mathrm{EM}\to e_5\to R_\psi\to e_4\to\alpha$
is broken at two independent steps (fixing $e_5$ and fixing $R_\psi$).
Therefore:

\[\boxed{\text{Alpha remains a free normalization in the current UBT formulation.}}\]
\end{theorem}

\paragraph{What is established:}
\begin{itemize}
  \item QED sector structure: DERIVED.
  \item U(1) generator normalization $\mathrm{Tr}(T^2)=1/2$: DERIVED.
  \item Integer charge quantization (unit flux convention): CONDITIONAL.
  \item RG running direction (decreasing $\alpha^{-1}$ with $\mu$): DERIVED.
  \item Qualitative consistency with checkpoints 137 and 127: OBSERVED CONSISTENCY.
\end{itemize}

\paragraph{What is not established:}
\begin{itemize}
  \item Absolute value of $e_5$: NO-GO.
  \item Absolute value of $R_\psi$: NO-GO.
  \item Numerical value of $\alpha$: NO-GO (follows from above).
\end{itemize}

\end{document}
